\section{积与图像}\label{sec:quasiprojective-products-and-graphs}

现在转向拟射影簇的积, 为研究正则映射的图像做准备. 图像将在下一节起核心作用. 为了说明拟射影情形下积的含义, 先回顾第~\ref{sec:products-and-segre-maps}~节的 Segre 映射
\[
  S_{m,n}:\mathbb{P}^{m}\times\mathbb{P}^{n}\to\mathbb{P}^{(m+1)(n+1)-1}.
\]
这是一个单射, 将一对射影簇的积 $X\times Y$ 送到 $\mathbb{P}^{(m+1)(n+1)-1}$ 中的一个射影簇. 这允许我们将 $X\times Y$ 与它在 $S_{m,n}$ 下的像等同, 从而赋予它射影簇结构. 基于大致相同的理由, 拟射影簇也有类似的结果.

\begin{Proposition}{拟射影簇的积}\leavevmode\label{prop:segre-image-of-quasiprojective-product}
若 $X\subseteq\mathbb{P}^{m}$ 和 $Y\subseteq\mathbb{P}^{n}$ 是拟射影簇, 则 Segre 映射 $S_{m,n}$ 限制为从 $X\times Y$ 到 $\mathbb{P}^{(m+1)(n+1)-1}$ 中的一个拟射影簇上的单射.
\end{Proposition}
\begin{Proof}
Proposition~\ref{prop:segre-map-injective-with-determinantal-image} 表明 $S_{m,n}$ 为单射, 所以只需证明 $S_{m,n}(X\times Y)$ 是拟射影簇. 为此, 首先利用 $X$ 为拟射影簇的假设, 写 $X=Z\setminus Z'$ 且 $Y=W\setminus W'$, 其中 $Z,Z'$ 在 $\mathbb{P}^{m}$ 中闭, $W,W'$ 在 $\mathbb{P}^{n}$ 中闭. 从集合论来看, 注意
\[
  X\times Y=(Z\times W)\setminus\bigl((Z\times W')\cup(Z'\times W)\bigr).
\]
由于 $S_{m,n}$ 为单射, 得到
\[
  S_{m,n}(X\times Y)=S_{m,n}(Z\times W)
   \setminus\bigl(S_{m,n}(Z\times W')\cup S_{m,n}(Z'\times W)\bigr).
\]
由 Proposition~\ref{prop:segre-map-restricts-to-product-of-projective-varieties}, 集合 $S_{m,n}(Z\times W)$, $S_{m,n}(Z\times W')$ 和 $S_{m,n}(Z'\times W)$ 都闭. 这表明 $S_{m,n}(X\times Y)$ 是一个闭集在另一个闭集中的补集, 换句话说, 它是拟射影簇.
\end{Proof}

据此命题, 可以将 $X\times Y$ 与 $S_{m,n}(X\times Y)$ 等同, 从而在拟射影簇的范畴内定义积的含义.

\begin{Definition}{拟射影簇的积}\leavevmode\label{def:product-of-quasiprojective-varieties}
若 $X\subseteq\mathbb{P}^{m}$ 和 $Y\subseteq\mathbb{P}^{n}$ 为拟射影簇, 则它们的\emph{积} $X\times Y$ (作为拟射影簇) 是拟射影簇
\[
  S_{m,n}(X\times Y)\subseteq\mathbb{P}^{(m+1)(n+1)-1}.
\]
\end{Definition}

应该停下来确认, 积的这一定义与我们先前在仿射和射影情形下的概念一致. 在射影情形下, 由定义立即可见它们一致. 在仿射情形下, 只须证明 $S_{m,n}$ 限制为从 $\mathbb{A}^{m+n}=\mathbb{A}^{m}_{0}\times\mathbb{A}^{n}_{0}$ 到其像上的 (拟射影簇的) 同构 (Exercise~\ref{ex:segre-map-restricts-to-affine-product-isomorphism}).

现在, 在拟射影簇的范畴内, 可以说明射影簇与仿射簇的积等对象的含义, 并且可以讨论所得积中的闭子集与开子集.

\begin{Example}{积 $\mathbb{P}^{1}\times\mathbb{A}^{2}$}\leavevmode\label{ex:product-of-projective-line-and-affine-plane}
考察位于 $\mathbb{P}^{1}\times\mathbb{P}^{2}$ 中的积 $\mathbb{P}^{1}\times\mathbb{A}^{2}$. 按定义, Segre 映射 $S_{1,2}:\mathbb{P}^{1}\times\mathbb{P}^{2}\to\mathbb{P}^{5}$ 由下式给出:
\[
  S_{1,2}([a_{0}:a_{1}],[b_{0}:b_{1}:b_{2}])
  =[a_{0}b_{0}:a_{0}b_{1}:a_{0}b_{2}:a_{1}b_{0}:a_{1}b_{1}:a_{1}b_{2}],
\]
相应地将 $\mathbb{P}^{5}$ 上的坐标记为 $[z_{00}:z_{01}:z_{02}:z_{10}:z_{11}:z_{12}]$. 有 $\mathbb{A}^{2}=\mathbb{A}^{2}_{0}=\mathbb{P}^{2}\setminus\mathcal{V}(y_{0})$, 读者可验证
\[
  S_{1,2}(\mathbb{P}^{1}\times\mathcal{V}(y_{0}))
    =S_{1,2}(\mathbb{P}^{1}\times\mathbb{P}^{2})\cap\mathcal{V}(z_{00},z_{10}).
\]
结合此式与 Proposition~\ref{prop:segre-map-injective-with-determinantal-image} 中 $S_{1,2}(\mathbb{P}^{1}\times\mathbb{P}^{2})$ 的表达式, 得到 $S_{1,2}(\mathbb{P}^{1}\times\mathbb{A}^{2})$ 是 $\mathbb{P}^{5}$ 中由下式给出的拟射影簇:
\begin{equation}\label{eq:segre-description-of-projective-line-times-affine-plane}
  \begin{aligned}
  &\mathcal{V}(z_{00}z_{11}-z_{10}z_{01},\ 
      z_{00}z_{12}-z_{10}z_{02},\ z_{01}z_{12}-z_{11}z_{02})\\
  &\qquad{}\setminus\mathcal{V}(z_{00},z_{10}).
  \end{aligned}
\end{equation}
依照 Definition~\ref{def:product-of-quasiprojective-varieties}, 将这个拟射影簇记为 $\mathbb{P}^{1}\times\mathbb{A}^{2}$. 那么, $\mathbb{P}^{1}\times\mathbb{A}^{2}$ 的一个闭子集的例子是集合
\[
  Z=\{([a_{0}:a_{1}],(b_{1},b_{2}))\mid a_{0}-a_{1}=0,\ b_{1}-1=0\},
\]
因为 $S_{1,2}$ 将 $Z$ 与 \eqref{eq:segre-description-of-projective-line-times-affine-plane} 同射影簇 $\mathcal{V}(z_{00}-z_{01},z_{01}-z_{10})$ 的交等同; 见 Exercise~\ref{ex:verify-closed-subset-of-projective-line-times-affine-plane}.
\end{Example}

若 $X\subseteq\mathbb{P}^{m}$ 和 $Y\subseteq\mathbb{P}^{n}$ 为拟射影簇, 则 $X\times Y$ 上的 Zariski 拓扑是从 $\mathbb{P}^{m}\times\mathbb{P}^{n}$ 继承的子空间拓扑. 如 Proposition~\ref{prop:projective-subvarieties-of-product-are-vanishing-sets} 所见, $\mathbb{P}^{m}\times\mathbb{P}^{n}$ 的闭集恰是 $K[x_{0},\ldots,x_{m},y_{0},\ldots,y_{n}]$ 中多项式族的消没集. 因此, $X\times Y$ 的闭子集对应于形如
\[
  (X\times Y)\cap\mathcal{V}(S)
\]
的交集, 其中 $S\subseteq K[x_{0},\ldots,x_{m},y_{0},\ldots,y_{n}]$.

为下一节的发展作准备, 现在转来讨论正则映射的图像. 这自然地推广了微积分中学习的实值函数图像的概念.

\begin{Definition}{正则映射的图像}\leavevmode\label{def:graph-of-quasiprojective-regular-map}
令 $F:X\to Y$ 为拟射影簇之间的正则映射. $F$ 的\emph{图像}是集合
\[
  \Gamma_{F}=\{(a,b)\in X\times Y\mid b=F(a)\}\subseteq X\times Y.
\]
\end{Definition}

例如, 若 $F:\mathbb{A}^{1}_{\mathbb{R}}\to\mathbb{A}^{1}_{\mathbb{R}}$ 是由 $F(a)=a^{2}$ 定义的多项式映射, 则图像 $\Gamma_{F}$ 就是 $\mathbb{A}^{1}_{\mathbb{R}}\times\mathbb{A}^{1}_{\mathbb{R}}=\mathbb{A}^{2}_{\mathbb{R}}$ 中通常的抛物线.

再来看一个略一般的例子.

\begin{Example}{一个正则映射的图像}\leavevmode\label{ex:graph-of-map-from-punctured-projective-line}
令 $X=\mathbb{P}^{1}\setminus\mathcal{V}(x_{0}x_{1})$, $Y=\mathbb{P}^{1}$. 则由
\[
  F([a_{0}:a_{1}])=[a_{0}a_{1}:a_{0}^{2}+a_{1}^{2}]
\]
给出的映射 $F:X\to Y$ 正则, 其图像可以描述为
\[
  \Gamma_{F}=\{([a_{0}:a_{1}],[b_{0}:b_{1}])\in X\times Y
    \mid[b_{0}:b_{1}]=[a_{0}a_{1}:a_{0}^{2}+a_{1}^{2}]\}.
\]
条件 $[b_{0}:b_{1}]=[a_{0}a_{1}:a_{0}^{2}+a_{1}^{2}]$ 等价于交叉乘积相等, 即 $b_{0}(a_{0}^{2}+a_{1}^{2})=b_{1}(a_{0}a_{1})$, 所以可以写
\[
  \Gamma_{F}=(X\times Y)\cap\mathcal{V}\bigl(y_{0}(x_{0}^{2}+x_{1}^{2})-y_{1}(x_{0}x_{1})\bigr)
    \subseteq\mathbb{P}^{1}\times\mathbb{P}^{1}.
\]
由于 $\Gamma_{F}$ 是 $X\times Y$ 与 $\mathbb{P}^{1}\times\mathbb{P}^{1}$ 中一个闭集的交, 可知 $\Gamma_{F}$ 是拟射影簇 $X\times Y$ 的闭子集.
\end{Example}

Example~\ref{ex:graph-of-map-from-punctured-projective-line} 末尾的观察说明了所有正则映射的图像都具有的一个性质. 下面的命题证明这一点.

\begin{Proposition}{正则映射的图像是闭的}\leavevmode\label{prop:graphs-of-quasiprojective-regular-maps-are-closed}
令 $X\subseteq\mathbb{P}^{m}$ 和 $Y\subseteq\mathbb{P}^{n}$ 为拟射影簇, $F:X\to Y$ 为正则映射. 则 $\Gamma_{F}$ 在 $X\times Y$ 中闭.
\end{Proposition}
\begin{Proof}
令 $(p,q)\in X\times Y$. 我们将构造 $(p,q)$ 的一个邻域 $U_{p}$ (只依赖于 $p$), 使 $\Gamma_{F}\cap U_{p}$ 在 $U_{p}$ 中闭. 由 Lemma~\ref{lem:closedness-is-local}, 可推出 $\Gamma_{F}$ 在 $X\times Y$ 中闭.

由于 $F$ 正则, 可选取同次数的齐次多项式 $f_{0},\ldots,f_{n}\in K[x_{0},\ldots,x_{m}]$, 使 $p\notin\mathcal{V}(f_{0},\ldots,f_{n})$, 且
\[
  F(a)=[f_{0}(a):\cdots:f_{n}(a)]
  \qquad\text{对所有 }a\in X\setminus\mathcal{V}(f_{0},\ldots,f_{n}).
\]
将 $f_{0},\ldots,f_{n}$ 视为 $K[x_{0},\ldots,x_{m},y_{0},\ldots,y_{n}]$ 中只依赖于 $x$ 变量的多项式, 定义开集
\[
  U_{p}=(X\times Y)\setminus\mathcal{V}(f_{0},\ldots,f_{n}),
\]
它包含点 $(p,q)$. 定义
\[
  Z_{p}=\mathcal{V}(\{y_{i}f_{j}-y_{j}f_{i}\mid0\le i,j\le n\})
    \subseteq\mathbb{P}^{m}\times\mathbb{P}^{n}.
\]
注意 $\Gamma_{F}\cap U_{p}=Z_{p}\cap U_{p}$, 因为对每个 $(a,b)\in U_{p}$, 条件 $b_{i}f_{j}(a)=b_{j}f_{i}(a)$ 对所有 $i,j$ 成立, 等价于要求
\[
  [b_{0}:\cdots:b_{n}]=[f_{0}(a):\cdots:f_{n}(a)]=F(a).
\]
由 Proposition~\ref{prop:projective-subvarieties-of-product-are-vanishing-sets}, $Z_{p}$ 在 $\mathbb{P}^{m}\times\mathbb{P}^{n}$ 中闭, 所以 $\Gamma_{F}\cap U_{p}=Z_{p}\cap U_{p}$ 在 $U_{p}$ 中闭, 如所需.
\end{Proof}

Proposition~\ref{prop:graphs-of-quasiprojective-regular-maps-are-closed} 对下文的重要性在于, 若 $F:X\to Y$ 为正则映射, 则可将像 $F(X)\subseteq Y$ 等同于 $\Gamma_{F}\subseteq X\times Y$ 在投影 $X\times Y\to Y$ 下的像. 因此, 理解 $\Gamma_{F}$ 并仔细研究投影映射, 将使我们能够对正则映射的像作出论断. 这就是本章最后一节的内容.

\BookExerciseHeading
\begin{BookExercises}
  \item\label{ex:segre-image-of-product-of-projective-open-sets-is-open} 证明若 $U\subseteq\mathbb{P}^{m}$ 和 $V\subseteq\mathbb{P}^{n}$ 为开集, 则像
  \[
    S_{m,n}(U\times V)\subseteq\mathbb{P}^{(m+1)(n+1)-1}
  \]
  在 $S_{m,n}(\mathbb{P}^{m}\times\mathbb{P}^{n})$ 中开.
  \item\label{ex:segre-map-restricts-to-affine-product-isomorphism}
  \begin{enumerate}
    \item 证明 Segre 映射 $S_{m,n}$ 限制为仿射空间 $\mathbb{A}^{m+n}=\mathbb{A}^{m}_{0}\times\mathbb{A}^{n}_{0}$ 与它在 $\mathbb{P}^{(m+1)(n+1)-1}$ 中的像之间的同构.
    \item 令 $X\subseteq\mathbb{A}^{m}$ 和 $Y\subseteq\mathbb{A}^{n}$ 为闭集. 证明闭集 $X\times Y\subseteq\mathbb{A}^{m+n}$ 与 Definition~\ref{def:product-of-quasiprojective-varieties} 中定义的拟射影簇 $X\times Y$ 同构. (提示: 使用 (a) 和 Corollary~\ref{cor:quasiprojective-isomorphisms-identify-subvarieties}.)
  \end{enumerate}
  \item\label{ex:verify-closed-subset-of-projective-line-times-affine-plane} 考察 Segre 映射 $S_{1,2}:\mathbb{P}^{1}\times\mathbb{P}^{2}\to\mathbb{P}^{5}$, 坐标的记号如 Example~\ref{ex:product-of-projective-line-and-affine-plane} 所示. 按以下步骤补充 Example~\ref{ex:product-of-projective-line-and-affine-plane} 的细节.
  \begin{enumerate}
    \item 验证 $S_{1,2}(\mathbb{P}^{1}\times\mathcal{V}(y_{0}))=S_{1,2}(\mathbb{P}^{1}\times\mathbb{P}^{2})\cap\mathcal{V}(z_{00},z_{10})$.
    \item 验证 Example~\ref{ex:product-of-projective-line-and-affine-plane} 中定义的集合 $Z$ 经 $S_{1,2}$ 映到集合 $S_{1,2}(\mathbb{P}^{1}\times\mathbb{A}^{2})\cap\mathcal{V}(z_{00}-z_{01},z_{01}-z_{10})$.
  \end{enumerate}
  \item\label{ex:product-of-closed-quasiprojective-subvarieties-is-closed} 令 $X_{1},X_{2}$ 为拟射影簇, $Y_{1}\subseteq X_{1}$ 和 $Y_{2}\subseteq X_{2}$ 为闭子簇. 证明 $Y_{1}\times Y_{2}$ 在 $X_{1}\times X_{2}$ 中闭.
  \item\label{ex:quasiprojective-product-projections-are-regular} 令 $X,Y$ 为拟射影簇. 证明投影映射 $\pi_{1}:X\times Y\to X$ 和 $\pi_{2}:X\times Y\to Y$ 正则.
  \item\label{ex:quasiprojective-products-of-regular-maps} 证明若 $F:Z\to X$ 与 $G:Z\to Y$ 是拟射影簇之间的正则映射, 则由
  \[
    (F\times G)(a)=(F(a),G(a))
  \]
  定义的映射 $F\times G:Z\to X\times Y$ 正则.
  \item\label{ex:quasiprojective-products-and-isomorphisms} 设 $X_{1},X_{2},Y_{1},Y_{2}$ 为拟射影簇, 使 $X_{1}\cong X_{2}$ 且 $Y_{1}\cong Y_{2}$. 证明
  \[
    X_{1}\times X_{2}\cong Y_{1}\times Y_{2}.
  \]
  \item\label{ex:quasiprojective-product-irreducible-iff-factors-are} 令 $X,Y$ 为拟射影簇. 证明 $X\times Y$ 不可约, 当且仅当 $X$ 和 $Y$ 不可约.
  \item\label{ex:graph-of-regular-map-is-isomorphic-to-domain} 令 $F:X\to Y$ 为拟射影簇之间的正则映射. 证明图像 $\Gamma_{F}\subseteq X\times Y$ 同构于 $X$.
\end{BookExercises}
